\title{The Tiling Method}
\author{Max von Hippel}
\date{\today}

\documentclass[12pt]{article}

\usepackage{usenixy}
\usepackage{
  amsmath,mathtools,amssymb,fullpage,
  amsthm,centernot,float,diagbox,adjustbox,
  graphicx,caption,subcaption,listings,
  enumerate,xcolor,hyperref,color,soul,
  stmaryrd,amsrefs,titling,tikz,
  mathrsfs,etoolbox,refcount,multicol,MnSymbol,
  yfonts,framed,relsize,colortbl,pgf-umlsd,varwidth,
  environ,xparse}
\lstset{
  language=Lisp,
  basicstyle=\ttfamily,
}
\usepackage[many]{tcolorbox}


\usepackage[framemethod=tikz]{mdframed}

\newcommand{\grayleftbar}[1]{\textcolor[gray]{0.7}{%
     \begin{leftbar}\textcolor{black}{#1}\end{leftbar}}}   
\newcommand{\darkleftbar}[1]{\textcolor[gray]{0.3}{%
     \begin{leftbar}\textcolor{black}{#1}\end{leftbar}}}   

\newcounter{countitems}
\newcounter{nextitemizecount}
\newcommand{\setupcountitems}{%
  \stepcounter{nextitemizecount}%
  \setcounter{countitems}{0}%
  \preto\item{\stepcounter{countitems}}%
}
\makeatletter
\newcommand{\computecountitems}{%
  \edef\@currentlabel{\number\c@countitems}%
  \label{countitems@\number\numexpr\value{nextitemizecount}-1\relax}%
}
\newcommand{\nextitemizecount}{%
  \getrefnumber{countitems@\number\c@nextitemizecount}%
}
\newcommand{\previtemizecount}{%
  \getrefnumber{countitems@\number\numexpr\value{nextitemizecount}-1\relax}%
}
\makeatother    
\newenvironment{AutoMultiColItemize}{%
\ifnumcomp{\nextitemizecount}{>}{3}{\begin{multicols}{2}}{}%
\setupcountitems\begin{itemize}}%
{\end{itemize}%
\unskip\computecountitems\ifnumcomp{\previtemizecount}{>}{3}{\end{multicols}}{}}

\usepackage{titlesec}
\usetikzlibrary{
    arrows,shapes,positioning,decorations.markings,
    decorations.pathreplacing,hobby,automata,tikzmark}
\tikzset{
    looped/.style={
        decoration={markings,mark=at position 0.999 with {\arrow[scale=2]{>}}},
        postaction={decorate},
        >=stealth
    },
    straight/.style={
        decoration={markings,mark=at position 1 with {\arrow[scale=2]{>}}},
        postaction={decorate},
        >=stealth
    },
    loopedSF/.style={
        decoration={
            markings,
            mark=at position 0.999 with {\arrow[scale=2]{>}},
            mark=at position 0.5 with {\arrow[scale=2]{>}}},
        postaction={decorate},
        >=stealth
    },
    straightSF/.style={
        decoration={
            markings,
            mark=at position 0.999 with {\arrow[scale=2]{>}},
            mark=at position 0.5 with {\arrow[scale=2]{>}}},
        postaction={decorate},
        >=stealth
    },
    triangle/.style = {fill=white, draw=black, regular polygon, regular polygon sides=3 },
    initial text=$$, % sets the text that appears on the start arrow
    node distance=1.5cm, % specifies the minimum distance between two nodes. Change if necessary.
    ->, % makes the edges directed
}
\titleformat{\subsection}[runin]
{\normalfont\large\bfseries}{\thesubsection}{1em}{}

\titleformat{\subsubsection}[runin]
{\normalfont\normalsize\bfseries}{\thesubsubsection}{1em}{}



\definecolor{mycolor}{rgb}{0.122, 0.435, 0.698}

\newmdenv[innerlinewidth=0.5pt, roundcorner=4pt,linecolor=mycolor,innerleftmargin=6pt,
innerrightmargin=6pt,innertopmargin=6pt,innerbottommargin=6pt]{mybox}

\usepackage[pdf]{graphviz}
\usetikzlibrary{shapes,arrows,calc}
\usepackage[utf8]{inputenc}

\setlength{\droptitle}{-7em}
\posttitle{\par\end{center}}

\newtheorem{definition}{Definition}
\newtheorem{theorem}{Theorem}
\newtheorem{lemma}{Lemma}
\newtheorem{problem}{Problem}
\newtheorem{example}{Example}
\newtheorem{remark}{Remark}

\newcommand{\TODO}[0]{\textcolor{red}{TODO}}
\newcommand{\F}[0]{\mathsf{F}}
\newcommand{\U}[0]{\mathsf{U}}
\newcommand{\G}[0]{\mathsf{G}}
\newcommand{\X}[0]{\mathsf{X}}
\newcommand{\acs}[0]{\textsc{ACL2S}}
\newcommand{\acl}[0]{\textsc{ACL2}}
\newcommand{\pr}[0]{\textsc{Promela}}
\newcommand{\ko}[0]{\textsc{Korg}}
\newcommand{\spn}[0]{\textsc{SPIN}}
\newcommand{\ie}{\emph{i.e.}}
\newcommand{\etal}{\emph{et al.}}

\newcommand{\Pcl}[0]{\textsc{P}}
\newcommand{\NP}[0]{\textsc{NP}}
\newcommand{\SAT}[0]{\texttt{SAT}}
\newcommand{\coNP}[0]{\textsc{coNP}}
\newcommand{\Time}[0]{\text{TIME}}
\newcommand{\Space}[0]{\text{SPACE}}
\newcommand{\threeSAT}[0]{\texttt{3SAT}}

\DeclarePairedDelimiter\abs{\lvert}{\rvert}
\DeclarePairedDelimiter\norm{\lVert}{\rVert}

\usepackage{tikz}
\usetikzlibrary{arrows,chains,matrix,positioning,scopes}

\newlength{\bubblewidth}
\AtBeginDocument{\setlength{\bubblewidth}{.75\textwidth}}
\definecolor{bubblegreen}{RGB}{103,184,104}
\definecolor{bubblegray}{RGB}{241,240,240}

\newcommand{\bubble}[4]{%
  \tcbox[
    colback=#1,
    colframe=#1,
    #2,
  ]{\color{#3}\begin{varwidth}{\bubblewidth}#4\end{varwidth}}%
}

\ExplSyntaxOn
\seq_new:N \l__ooker_bubbles_seq
\tl_new:N \l__ooker_bubbles_first_tl
\tl_new:N \l__ooker_bubbles_last_tl

\NewEnviron{rightbubbles}
 {
  \raggedleft\sffamily
  \seq_set_split:NnV \l__ooker_bubbles_seq { \par } \BODY
  \int_compare:nTF { \seq_count:N \l__ooker_bubbles_seq < 2 }
   {
    \bubble{bubblegreen}{rounded~corners}{white}{\BODY}
   }
   {
    \seq_pop_left:NN \l__ooker_bubbles_seq \l__ooker_bubbles_first_tl
    \seq_pop_right:NN \l__ooker_bubbles_seq \l__ooker_bubbles_last_tl
    \bubble{bubblegreen}{sharp~corners=southeast}{white}{\l__ooker_bubbles_first_tl}\par
    \seq_map_inline:Nn \l__ooker_bubbles_seq
     {
      \bubble{bubblegreen}{sharp~corners=east}{white}{##1}\par
     }
    \bubble{bubblegreen}{sharp~corners=northeast}{white}{\l__ooker_bubbles_last_tl}\par
   }
 }
\NewEnviron{leftbubbles}
 {
  \raggedright\sffamily
  \seq_set_split:NnV \l__ooker_bubbles_seq { \par } \BODY
  \int_compare:nTF { \seq_count:N \l__ooker_bubbles_seq < 2 }
   {
    \bubble{bubblegray}{rounded~corners}{black}{\BODY}
   }
   {
    \seq_pop_left:NN \l__ooker_bubbles_seq \l__ooker_bubbles_first_tl
    \seq_pop_right:NN \l__ooker_bubbles_seq \l__ooker_bubbles_last_tl
    \bubble{bubblegray}{sharp~corners=southwest}{black}{\l__ooker_bubbles_first_tl}\par
    \seq_map_inline:Nn \l__ooker_bubbles_seq
     {
      \bubble{bubblegray}{sharp~corners=west}{black}{##1}\par
     }
    \bubble{bubblegray}{sharp~corners=northwest}{black}{\l__ooker_bubbles_last_tl}\par
   }
 }
\ExplSyntaxOff

\begin{document}
\maketitle

Consider a two-party communication problem,
	in which the participants

\begin{figure}[h]
\centering
\begin{subfigure}{.3\textwidth}
  \centering
  \includegraphics[width=.2\linewidth]{alice.png}
  \caption{Alice}
  \label{fig:alice}
\end{subfigure}%
and
\begin{subfigure}{.3\textwidth}
  \centering
  \includegraphics[width=.2\linewidth]{bob.png}
  \caption{Bob}
  \label{fig:bob}
\end{subfigure}
\label{fig:participants}
\end{figure}

\emph{participate} to compute a function:
\[f : \underbrace{\mathbb{B}^n}_{\begin{aligned}\text{Alice's}\\\text{input}\end{aligned}} \times 
      \underbrace{\mathbb{B}^n}_{\begin{aligned}\text{Bob's}\\\text{input}\end{aligned}} \to 
      \underbrace{\mathbb{B}}_{\begin{aligned}\text{global}\\\text{output}\end{aligned}}\]
The players can come up with a \emph{protocol} 
\(\Pi = (p_1, ..., p_t)\),
namely, 
	for some natural $t \in \mathbb{N}$,
	a sequence of $t$-many functions
\(p_i : \mathbb{B}^* \to \mathbb{B}^*\)
such that the communication between the players looks like this:\\
\noindent\hfill\par
\noindent\makebox[\textwidth][c]{%
    \begin{minipage}{0.8\textwidth}
\noindent\fbox{\parbox{\textwidth}{
Alice is given input $x$.
\begin{leftbubbles}
Hi Bob. I'm not divulging $x$, but, $p_1(x) = \texttt{P1}$.
\end{leftbubbles}
\rightline{Bob is given input $y$.}
\begin{rightbubbles}
Thanks Alice.  I'm not divulging $y$, but, $p_2(y, \texttt{P1}) = \texttt{P2}$.
\end{rightbubbles}
\begin{leftbubbles}
Thanks Bob.  Don't tell anyone, but: $p_3(x, \texttt{P1}, \texttt{P2}) = \texttt{P3}$.
\end{leftbubbles}
\begin{rightbubbles}
Is that so? Well, $p_4(y, \texttt{P1}, \texttt{P2}, \texttt{P3}) = \texttt{P4}$.
\end{rightbubbles}
\begin{leftbubbles}
Wicked.  In that case, $p_5(x, \texttt{P1}, \texttt{P2}, \texttt{P3}, \texttt{P4}) = \texttt{P5}$.
\end{leftbubbles}
\[... \emph{yada yada yada} ...\]
\begin{leftbubbles}
In that case, the last thing you need to know is that $p_t(x, \texttt{P1}, \texttt{P2}, \texttt{P3}, \texttt{P4}, ..., \texttt{P(t-1)})=\texttt{Pt}$.
\end{leftbubbles}}}
\end{minipage}}
\clearpage
Suppose that there is a protocol $\Pi$ for $f$ consisting of $t$ messages,
but, there does not exist any protocol $\Pi'$ for $f$ consisting of fewer than $t$ messages.
Then we say $t$ is the \emph{communication complexity} of $f$, and we write $C(f) = t$.

Given some such function $f$, it would be nice if we could automatically compute
a reasonable lower bound on its communication complexity.
One way to do this is with the \emph{tiling method}.
We will give the method immediately, and in tandem, we will illustrate the method using
the function $f(x, y) = x < y$ where $x, y$ are integers in $\{ 0, 1, 2, 3 \}$, encoded in $\mathbb{B}$oolean.
First, let $M(f)$ be the \emph{matrix of }$f$,
	namely,
		the $2^n \times 2^n$ matrix whose $(x, y)$th entry
		is the value $f(x, y)$.

\begin{table}[h]
\centering
\begin{tabular}{|l|lllll|}
\hline
    & 0 & 1 & 2 & 3 & 4 \\\hline
\texttt{000} = 0   & \cellcolor{red!25}0 & \cellcolor{green!25}1 & \cellcolor{green!25}1 & \cellcolor{green!25}1 & \cellcolor{green!25}1 \\
\texttt{001} = 1   & \cellcolor{red!25}0 & \cellcolor{red!25}0 & \cellcolor{green!25}1 & \cellcolor{green!25}1 & \cellcolor{green!25}1 \\
\texttt{010} = 2   & \cellcolor{red!25}0 & \cellcolor{red!25}0 & \cellcolor{red!25}0 & \cellcolor{green!25}1 & \cellcolor{green!25}1 \\
\texttt{011} = 3   & \cellcolor{red!25}0 & \cellcolor{red!25}0 & \cellcolor{red!25}0 & \cellcolor{red!25}0 & \cellcolor{green!25}1 \\
\texttt{100} = 4   & \cellcolor{red!25}0 & \cellcolor{red!25}0 & \cellcolor{red!25}0 & \cellcolor{red!25}0 & \cellcolor{red!25}0\\\hline
\end{tabular}
\caption{The matrix $M(<)$ for inputs $x, y \in \{ 0, 1, 2, 3 \}$.  Values of $x$ are given in the rows, while values of $y$ are given in the columns.  False (i.e. 0) values are marked red for clarity.}
\end{table}

A \emph{combinatorial rectangle} in $M(f)$ is any submatrix of $M$.
We say a rectangle $A \times B$ in $M(f)$ is \emph{monochromatic}
if for all $x, x'$ in $A$ and $y, y'$ in $B$, $M_{x,y} = M_{x',y'}$.

\begin{figure}[h]
\centering 
	\begin{minipage}{0.25\textwidth}
		\begin{tabular}{|l|lllll|}
		\hline
		    & 0 & 1 & 2 & 3 & 4 \\\hline
		0   & 0 & 1 & 1 & 1 & 1 \\
		1   & 0 & 0 & 1 & 1 & 1 \\
		2   & \cellcolor{purple!75}{\textcolor{white}{0}} & \cellcolor{purple!75}{\textcolor{white}{0}} & \cellcolor{purple!75}{\textcolor{white}{0}} & 1 & 1 \\
		3   & \cellcolor{purple!75}{\textcolor{white}{0}} & \cellcolor{purple!75}{\textcolor{white}{0}} & \cellcolor{purple!75}{\textcolor{white}{0}} & 0 & 1 \\
		4   & \cellcolor{purple!75}{\textcolor{white}{0}} & \cellcolor{purple!75}{\textcolor{white}{0}} & \cellcolor{purple!75}{\textcolor{white}{0}} & 0 & 0\\\hline
		\end{tabular}
	\end{minipage}
	\begin{minipage}{0.25\textwidth}
		\begin{tabular}{|l|lllll|}
		\hline
		    & 0 & 1 & 2 & 3 & 4 \\\hline
		0   & \cellcolor{orange!75}0 & 1 & \cellcolor{orange!75}1 & 1 & \cellcolor{orange!75}1 \\
		1   & \cellcolor{orange!75}0 & 0 & \cellcolor{orange!75}1 & 1 & \cellcolor{orange!75}1 \\
		2   & 0 & 0 & 0 & 1 & 1 \\
		3   & \cellcolor{orange!75}0 & 0 & \cellcolor{orange!75}0 & 0 & \cellcolor{orange!75}1 \\
		4   & 0 & 0 & 0 & 0 & 0\\\hline
		\end{tabular}   
	\end{minipage}
	\begin{minipage}{0.25\textwidth}
		\begin{tabular}{|l|lllll|}
		\hline
		    & 0 & 1 & 2 & 3 & 4 \\\hline
		0   & 0 & 1 & \cellcolor{black!75}{\textcolor{white}{1}} & \cellcolor{black!75}{\textcolor{white}{1}} & 1 \\
		1   & 0 & 0 & \cellcolor{black!75}{\textcolor{white}{1}} & \cellcolor{black!75}{\textcolor{white}{1}} & 1 \\
		2   & 0 & 0 & \cellcolor{black!75}{\textcolor{white}{0}} & \cellcolor{black!75}{\textcolor{white}{1}} & 1 \\
		3   & 0 & 0 & 0 & 0 & 1 \\
		4   & 0 & 0 & 0 & 0 & 0\\\hline
		\end{tabular}   
	\end{minipage}
	\caption{Some example rectangles of $M(<)$.  The first rectange, in {\colorbox{purple!75}{\textcolor{white}{purple}}}, is monochromatically colored 0.  The second rectangle, in {\colorbox{orange!75}{orange}}, is illustrates the flexibility of our rectangle definition, namely, that the rectangle does not actually need to be connected in the original matrix (although, the entries cannot be permuted).  Neither the {\colorbox{purple!75}{\textcolor{white}{second}}} nor {\colorbox{black!75}{\textcolor{white}{third}}} rectangle is monochromatic.}
\end{figure}

Without loss of generality, 
	suppose the protocol $\Pi$ begins with Alice
	sending a bit $b$.
Then \emph{certainly} $M(f)$ partitions into two rectangles,
	where the first considers all the scenarios where the bit Alice sent was 0,
	and the second considers all the scenarios where the bit Alice sent was 1.
By induction, every bit sent in the protocol further \emph{naturally partitions} $M(f)$,
and we end up with a binary tree of partitions of $M(f)$ induced by $\Pi$.

% \bibliographystyle{IEEEtran}
% \bibliography{htm}
\end{document}